\documentclass[a4paper,12pt]{article}
\usepackage[utf8]{inputenc}
\usepackage[ngerman]{babel}
\usepackage{amsmath}
\usepackage{geometry}
\geometry{a4paper, top=15mm, left=15mm, right=15mm, bottom=15mm, 
	headsep=10mm, footskip=12mm}
\usepackage{parskip}
\usepackage{booktabs}
\usepackage{tikz}
\usepackage{graphicx}
\usepackage{href-ul}
\hypersetup{ 
	colorlinks=true, 
	linkcolor=blue, 
	urlcolor=blue}

\begin{document}
	
	

\begin{tikzpicture}(10,3)
	\draw[ultra thick](2,0) --(10,0) -- (10,1.3) --(2,1.3) -- (2,0);
	\draw[fill=black](2,0)-- (0,.6) -- (2,.6) -- (2,0);
	\node at (6,.6) {\large Solution Sheet of \href{https://www.okuyakl.de}{www.okuyakl.de}};
\end{tikzpicture}

	
	\section*{1. Different power sources – typical values}
	
	\begin{center}
		\begin{tabular}{ccc}
			\toprule
			\textbf{Power source} & \textbf{Voltage (V)} & \textbf{Example application} \\
			\midrule
			AA battery & 1.5 & torch, remote control \\
			Power bank (USB) & 5.0 & cell phone charging \\
			E-bike battery & 36 – 48 & electric motor on a bicycle \\
			Household socket & 230 & vacuum cleaner, lamps \\
			\bottomrule
		\end{tabular}
	\end{center}
	
	\hrule
	
	\section*{2. The 6-volt light bulb – what works and what doesn't?}
		
		\textbf{AA battery (1.5 V)}\\
		a) The lamp barely glows or doesn't glow at all – voltage too low.\\
		b) Not dangerous, but pointless.
		
		\textbf{Power bank (5 V USB)}\\
		a) The lamp glows, but slightly dimmer than at rated voltage.\\
		b) Rather useful if nothing else is available – not critical.
		
		\textbf{2 AA batteries in series (3 V)}\\
		a) The lamp glows slightly.\\
		b) Not dangerous, but too dim.
		
		\textbf{4 AA batteries in series (6 V)}\\
		a) The lamp glows perfectly.\\
		b) Perfectly suitable – full rated voltage.
		
		\textbf{E-bike battery (36 V)} \\
		a) The lamp burns out immediately. \\
		b) Extremely dangerous – not suitable!
		
		\textbf{Socket (230 V)} \\
		a) The lamp explodes or almost vaporizes. \\
		b) Life-threatening – do not do this under any circumstances!
		
		\hrule
		
		\section*{3. Experiment}
		
		\renewcommand{\arraystretch}{2}
		\begin{tabular}{|c|c|c|c|}
			\hline
			Voltage U & Power P & Operating Current I & Resistance R \\
			\hline
			6 V & 0.6 W & 0.1 A & 60 $\Omega$\\
			\hline
			6 V & 1.8 W & 0.3 A & 20 $\Omega$\\
			\hline
			3.5 V & 0.7 W & 0.2 A & 17.5 $\Omega$\\
			\hline
			2.2 V & 0.5 W & 0.23 A & 9.7 $\Omega$ \\
			\hline
		\end{tabular}
		\hrule
		
		\newpage
		
		Circuits for \textbf{one} light bulb:
		
		\textbf{Circuit 1:}
		Power bank with USB output (5 V) and 6 V lamp
		\textit{Lamp lights up, but not quite at full power}
		
		\textbf{Circuit 2:}
		One 1.5 V AA battery and 2.2 V lamp
		\textit{Lamp lights up, but not quite at full power}
		
		\textbf{Circuit 3:}
		Two AA batteries in series: 1.5 V times 2 = 3 V and 3.5 V lamp
		\textit{Lamp lights up, but not quite at full power}
		
		\textbf{Circuit 4:}
		Four AA batteries in series: 1.5 V times 4 = 6 V and 6 V lamp
		\textit{Ideally matched voltage, lamp lights up at full power}
		
		Circuits for \textbf{multiple} light bulbs:
		
		\textbf{Circuit 5:}
		Power bank V with USB output (5 V) and two 6 V bulbs in parallel
		\textit{Lamps light up, but not quite at full power}
		
		\textbf{Circuit 6:}
		Four AA batteries in series, 1.5 V times 4 = 6 V, and two 6 V bulbs in parallel
		\textit{Ideally matched voltage, bulbs light up at full power}
		
		\textbf{Circuit 7:}
		Eight AA batteries in series, 1.5 V times 4 = 12 V, and two identical 6 V bulbs in series
		\textit{Ideally matched voltage, bulbs light up at full power}
		
		\textbf{Circuit 8:}
		Four AA batteries in series, 1.5 V times 4 = 6 V, and three 2.2 V bulbs in series
		\textit{Lamps light up, but not at full power}
	
	\textbf{Circuit 9:} 
	Eight AA batteries in series 1.5 V times 4 = 12 V and four 3.5 V lamps in series \\
	\textit{Lamps light up, but not quite at full power}
	
	\begin{center}
		\includegraphics[width=7 cm]{../../viecher/eendcomic.pdf}
		
		Click here to return to the \href{https://www.okuyakl.de/phys/p7esdL131/ae131.pdf}{worksheet}
	\end{center}
\end{document}